(1) 设第 n 原胞内两个质点位移为 \(u_n,v_n\)，方程为
\[
m_1\ddot u_n=K(v_n+v_{n-1}-2u_n),\quad
m_2\ddot v_n=K(u_n+u_{n+1}-2v_n).
\]
弹簧回复力只依赖相邻质点的相对位移。

(2) 取 \(u_n=Ue^{i(nqa-\omega t)},v_n=Ve^{i(nqa-\omega t)}\)，构成
\[
(2K-m_1\omega^2)U-K(1+e^{-iqa})V=0,\quad
-K(1+e^{iqa})U+(2K-m_2\omega^2)V=0.
\]
行列式条件为 \(m_1m_2\omega^4-2K(m_1+m_2)\omega^2+4K^2\sin^2(qa/2)=0\)。解得
\[
\omega_\pm^2=K\left(\frac1{m_1}+\frac1{m_2}\right)
\pm K\sqrt{\left(\frac1{m_1}+\frac1{m_2}\right)^2-\frac{4\sin^2(qa/2)}{m_1m_2}}.
\]

(3) q=0 时 \(\omega_-(0)=0,\ \omega_+^2(0)=2K(1/m_1+1/m_2)\)。区边 \(|q|=\pi/a\) 时根退化为 \(\omega^2=2K/m_1\) 和 \(2K/m_2\)，若两质量不等，区边两个分支之间存在频率间隔。

(4) 负号支在 q→0 时 \(\omega_-\propto|q|\)，两质点主要同相，称声学支；正号支在 q=0 仍有非零频率，原胞内部相对振动，称光学支。当 \(m_1=m_2=m\)，区边分支相接于 \(\omega^2=2K/m\)，频隙关闭，两个分支是单原子链在双倍原胞下的能带折叠。
